\documentclass[11pt,a4paper]{article}
\usepackage[utf8]{inputenc}
\usepackage[T1]{fontenc}
\usepackage[portuguese]{babel}
\usepackage{xcolor}
\usepackage{hyperref}
\usepackage{geometry}
\usepackage{listings}
\usepackage{tcolorbox}
\usepackage{graphicx}

\geometry{margin=2.5cm}

\definecolor{primary}{RGB}{41, 128, 185}
\definecolor{secondary}{RGB}{44, 62, 80}
\definecolor{codebg}{RGB}{240, 240, 240}
\definecolor{textcol}{RGB}{51, 51, 51}

\hypersetup{
    colorlinks=true,
    linkcolor=primary,
    filecolor=magenta,      
    urlcolor=primary,
}

\lstset{
    backgroundcolor=\color{codebg},
    basicstyle=\ttfamily\small\color{secondary},
    breaklines=true,
    frame=single,
    rulecolor=\color{codebg},
    keywordstyle=\color{primary}\bfseries,
    commentstyle=\color{gray}\itshape,
    stringstyle=\color{teal},
    showstringspaces=false
}

\title{\color{secondary}\Huge\textbf{Componentes Internos do Kubernetes}}
\author{\color{primary}DevOps com Kubernetes -- 2026}
\date{}

\begin{document}
\maketitle

\section*{\color{primary}O que você aprenderá nesta página}
\begin{itemize}
    \item Explicar, em termos gerais, como o Kubernetes opera.
\end{itemize}

Em vez de pensar no Kubernetes como algo completamente novo, percebi que compará-lo a um sistema operacional ajuda bastante. Não sou nenhum especialista em sistemas operacionais, mas todos nós os usamos.

O Kubernetes é uma camada sobre a qual executamos nossas aplicações. Ele pega os recursos que estão acessíveis a partir das camadas inferiores e gerencia nossas aplicações e recursos. Ele também fornece serviços, como DNS, para as aplicações. Com essa mentalidade de sistema operacional, também podemos tentar fazer o caminho inverso: você talvez já tenha utilizado um cron (ou o \textit{task scheduler} no Windows) para agendar tarefas em lote (\textit{batch jobs}), como salvar os backups de algumas aplicações. A mesma coisa pode ser feita no Kubernetes com CronJobs.

Agora que começaremos a falar sobre o funcionamento interno, obteremos novos conhecimentos a respeito do Kubernetes e seremos capazes de prever e resolver problemas que decorram de sua natureza.

Como esta seção é principalmente uma reiteração da documentação do Kubernetes, incluirei vários links para a documentação oficial. Apesar de estarmos falando sobre os componentes internos, não discutiremos como montar seu próprio cluster do zero manualmente. Se você deseja "sujar as mãos" e configurá-lo, deveria ler e completar Kubernetes the Hard Way, escrito por Kelsey Hightower.

\subsection*{\color{secondary}Controladores (\textit{Controllers})}
Os Controladores monitoram o estado do seu cluster e tentam sempre aproximar o estado atual do estado desejado. Quando você declara \textit{X} réplicas de um Pod no seu \texttt{deployment.yaml}, um controlador chamado Replication Controller garantirá que isso seja mantido como verdade. Existem controladores específicos para várias responsabilidades.

\subsection*{\color{secondary}Plano de Controle do Kubernetes (\textit{Kubernetes Control Plane})}
O Plano de Controle do Kubernetes (\textit{Control Plane}) é responsável pelo gerenciamento de todo o cluster do Kubernetes. É o principal componente de orquestração responsável por garantir que o estado desejado do cluster seja correspondente ao estado atual. O Control Plane toma decisões globais sobre o cluster (como, por exemplo, agendamento de tarefas) e também detecta e responde aos eventos do cluster (como o início de um novo Pod caso o número de réplicas de um deployment não esteja sendo satisfeito).

O Plano de Controle é formado por:
\begin{itemize}
    \item \textbf{etcd}: Um armazenamento chave-valor (\textit{key-value storage}) que o Kubernetes usa para salvar e respaldar todos os dados do cluster.
    \item \textbf{kube-scheduler}: Decide em qual \textit{node} (nó) um Pod deverá ser executado.
    \item \textbf{kube-controller-manager}: É o responsável por executar e controlar todos os controladores.
    \item \textbf{kube-apiserver}: Expõe o \textit{Control Plane} do Kubernetes aos usuários por meio de uma API.
\end{itemize}

Existe também um componente extra chamado \textit{cloud-controller-manager}, que permite vincular o cluster à API de um provedor de Nuvem. Se você quisesse montar seu próprio cluster na Hetzner, por exemplo, poderia utilizar o \texttt{hcloud-cloud-controller-manager} dentro do seu próprio cluster Kubernetes sendo hospedado através de máquinas virtuais (VMs) deles.

\subsection*{\color{secondary}Componentes de Nó (\textit{Node Components})}
Cada nó (node) possui alguns componentes com o intuito de manter os \textit{Pods} funcionando:
\begin{itemize}
    \item \textbf{kubelet}: Garante que os contêineres estão executando as suas rotinas em um \textit{Pod}.
    \item \textbf{kube-proxy}: É um proxy e um mantenedor de regras de rede. Permite conexões com a parte interna e externa do cluster, para que os Serviços operem conforme os usamos.
\end{itemize}
Além disso, cada nó tem o chamado \textit{Container Runtime}. Neste curso, estamos usando o \textbf{Docker} como tempo de execução.

\subsection*{\color{secondary}Componentes Adicionais (\textit{Addons})}
Apesar de tudo o que foi citado acima, o Kubernetes permite componentes adicionais (\textit{Addons}), que utilizam os mesmos recursos do Kubernetes já discutidos e estendem suas funcionalidades. Você pode inspecionar quais recursos foram criados por essas extensões usando o comando no \textit{namespace} \texttt{kube-system}:

\begin{lstlisting}[language=bash]
$ kubectl -n kube-system get all
NAME                                          READY   STATUS      RESTARTS   AGE
pod/coredns-ccb96694c-t9rr2                   1/1     Running     0          42h
pod/helm-install-traefik-crd-5446f            0/1     Completed   0          42h
pod/helm-install-traefik-p8c2h                0/1     Completed   2          42h
pod/local-path-provisioner-5cf85fd84d-s9x7p   1/1     Running     0          42h
pod/metrics-server-5985cbc9d7-hngw9           1/1     Running     0          42h
pod/svclb-traefik-68897c6e-bq9k7              2/2     Running     0          42h
pod/svclb-traefik-68897c6e-kw5p4              2/2     Running     0          42h
pod/svclb-traefik-68897c6e-qkc5c              2/2     Running     0          42h
pod/traefik-5d45fc8cc9-ws66x                  1/1     Running     0          42h

NAME                     TYPE           CLUSTER-IP     EXTERNAL-IP                                 PORT(S)                      AGE
service/kube-dns         ClusterIP      10.43.0.10     <none>                                      53/UDP,53/TCP,9153/TCP       42h
service/metrics-server   ClusterIP      10.43.224.88   <none>                                      443/TCP                      42h
service/traefik          LoadBalancer   10.43.140.57   192.168.144.3,192.168.144.4,192.168.144.5   80:32558/TCP,443:32155/TCP   42h
\end{lstlisting}

Para ter uma visão abrangente sobre como cada parte conversa entre si, leia o artigo \textit{what happens when k8s}, que explora o que ocorre de forma interna quando você executa um comando como \texttt{kubectl run nginx --image=nginx --replicas=3}, o que ilumina de verdade a "mágica" que atua por debaixo dos panos.

\section*{\color{primary}Auto-recuperação (\textit{Self-healing})}

Anteriormente, conversamos brevemente a respeito da natureza de auto-recuperação do Kubernetes e sobre como \textit{Pods} que são deletados podem ser recriados de maneira automática.

Mas o que ocorre se apagarmos ou derrubarmos um nó (node) contendo um \textit{pod}? Primeiramente, faremos o \textit{deploy} de um aplicativo web com \textit{ingress}, e confirmaremos onde seu \textit{pod} está sendo executado:

\begin{lstlisting}[language=bash]
$ kubectl apply -f https://raw.githubusercontent.com/kubernetes-hy/material-example/master/app2/manifests/deployment.yaml \
                -f https://raw.githubusercontent.com/kubernetes-hy/material-example/master/app2/manifests/ingress.yaml \
                -f https://raw.githubusercontent.com/kubernetes-hy/material-example/master/app2/manifests/service.yaml
  deployment.apps/hashresponse-dep created
  ingress.extensions/dwk-material-ingress created
  service/hashresponse-svc created

$ curl localhost:8081
9eaxf3: 8k2deb

$ kubectl describe po hashresponse-dep-57bcc888d7-5gkc9 | grep 'Node:'
  Node:         k3d-k3s-default-agent-1/172.30.0.2
\end{lstlisting}

Vemos aqui que ele está localizado em \texttt{agent-1}. Em seguida, faremos com que este nó saia da rede, tirando-o do ar usando um "pause" via docker:

\begin{lstlisting}[language=bash]
$ docker ps
  CONTAINER ID        IMAGE                      COMMAND                  CREATED             STATUS              PORTS                                           NAMES
  5c43fe0a936e        rancher/k3d-proxy:v3.0.0   "/bin/sh -c nginx-pr..." 10 days ago         Up 2 hours          0.0.0.0:8081->80/tcp, 0.0.0.0:50207->6443/tcp   k3d-k3s-default-serverlb
  fea775395132        rancher/k3s:latest         "/bin/k3s agent"         10 days ago         Up 2 hours                                                          k3d-k3s-default-agent-1
  70b68b040360        rancher/k3s:latest         "/bin/k3s agent"         10 days ago         Up 2 hours          0.0.0.0:8082->30080/tcp                         k3d-k3s-default-agent-0
  28cc6cef76ee        rancher/k3s:latest         "/bin/k3s server --t..." 10 days ago         Up 2 hours                                                          k3d-k3s-default-server-0

$ docker pause k3d-k3s-default-agent-1
k3d-k3s-default-agent-1
\end{lstlisting}

Aguardaremos alguns minutos, e este deverá ser o estado que encontraremos:

\begin{lstlisting}[language=bash]
$ kubectl get po
NAME                                READY   STATUS        RESTARTS   AGE
hashresponse-dep-57bcc888d7-5gkc9   1/1     Terminating   0          15m
hashresponse-dep-57bcc888d7-4klvg   1/1     Running       0          30s

$ curl localhost:8081
6xluh2: ta0ztp
\end{lstlisting}

E caso seja o único \textit{control-plane} (Plano de Controle) que acabe deletado ou offline? Não ocorrerá nada de bom. Sendo um cluster local, ele se mostra um ponto de falha única (SPOF - Single Point of Failure). Confira a documentação do Kubernetes sobre Options for Highly Available topology para evitar que seu cluster inteiro fique inválido devido a uma simples falha de hardware.

\end{document}
